\chapter{HIV and AIDS Medicines}

\section*{Definition and Overview}
HIV medicines, known as antiretroviral therapy (ART), are combinations of drugs that stop the human immunodeficiency virus from multiplying. They work by blocking different stages of the viral life cycle. ART does not cure HIV, but it reduces the amount of virus in the blood to undetectable levels, which allows the immune system to recover, prevents progression to AIDS, and makes transmission impossible. Modern treatment usually consists of one or two tablets taken daily, is well tolerated, and is recommended for everyone living with HIV from the time of diagnosis.

\section*{Epidemiology}
Around 30,000 Australians live with HIV, and the vast majority now take antiretroviral therapy. Treatment has transformed HIV from a fatal illness to a manageable chronic condition, and people on effective treatment have near-normal life expectancy. In Australia, over 90\% of people diagnosed with HIV are on treatment, and most have an undetectable viral load. The success of treatment has contributed to falling rates of new HIV diagnoses and near-elimination of mother-to-child transmission.

\section*{Aetiology and Pathophysiology}
\begin{itemize}
    \item \textbf{Classes of medicines:} nucleoside reverse transcriptase inhibitors, non-nucleoside reverse transcriptase inhibitors, protease inhibitors, integrase inhibitors, and entry inhibitors
    \item \textbf{Mechanism:} each class blocks a different step in the viral life cycle, so the virus cannot replicate
    \item \textbf{Combination therapy:} three drugs from at least two classes are used to prevent resistance
    \item \textbf{Integrase inhibitors:} now the preferred backbone, highly potent with few side effects
    \item \textbf{Pathophysiology:} without replication, viral load falls to undetectable, CD4 cells recover, and the immune system rebuilds
    \item \textbf{Resistance:} occurs when treatment is missed or viral replication continues
    \item \textbf{Modern regimens:} many people take a single daily tablet containing several drugs
\end{itemize}

\section*{Clinical Features}
\begin{itemize}
    \item Treatment is started as soon as possible after diagnosis
    \item Viral load becomes undetectable within 3--6 months in nearly all people
    \item CD4 count rises as immune function recovers
    \item Most people feel well and experience few or no side effects
    \item Some people have initial side effects such as nausea or headache that settle
    \item Long-term use is monitored for effects on kidneys, bones, lipids, and metabolism
    \item Medicines are adjusted if resistance or side effects occur
\end{itemize}

\section*{Differential Diagnosis}
\begin{itemize}
    \item Drug interactions with other medicines are common and must be reviewed
    \item Side effects can be confused with symptoms of other conditions
    \item Resistance requires virological testing to distinguish from adherence problems
    \item Co-infections such as hepatitis B and C need their own treatment
    \item Regimens are individualised based on resistance testing, pregnancy, kidney function, and other factors
\end{itemize}

\section*{When to Seek Medical Care}
Anyone newly diagnosed with HIV should start antiretroviral therapy promptly through specialist or GP care. People on treatment should attend regular reviews and report new symptoms, side effects, or difficulty taking medicines.

\textbf{Call 000 (or local emergency services) for:}
\begin{itemize}
    \item Severe allergic reaction to a medicine: swelling, rash, or difficulty breathing
    \item Severe abdominal pain, vomiting, or jaundice (possible liver toxicity)
    \item Confusion, seizures, or collapse
    \item Signs of a serious infection with fever and drowsiness
    \item Thoughts of self-harm or suicide
\end{itemize}

\section*{Diagnosis and Investigations}
\begin{itemize}
    \item \textbf{Viral load testing:} monitors treatment response; undetectable is the goal
    \item \textbf{CD4 count:} tracks immune recovery
    \item \textbf{Drug resistance testing:} before starting or changing treatment
    \item \textbf{Blood tests:} kidney and liver function, blood count, lipids, and blood sugar
    \item \textbf{Hepatitis and tuberculosis screening:} because of co-infections
    \item \textbf{Pregnancy testing:} relevant in women planning pregnancy
    \item \textbf{Regular monitoring:} every 3--6 months while stable
\end{itemize}

\section*{Management}
\begin{itemize}
    \item \textbf{Standard treatment:} a single daily tablet combining three drugs, usually including an integrase inhibitor
    \item \textbf{Long-acting injectables:} monthly or two-monthly injections are available for some people
    \item \textbf{Adherence:} near-perfect adherence is needed to maintain suppression
    \item \textbf{Treatment in pregnancy:} safe regimens prevent transmission to the baby
    \item \textbf{Switching:} regimens are changed for side effects, interactions, or resistance
    \item \textbf{Prevention of transmission:} undetectable viral load means no transmission
    \item \textbf{Management of coexisting conditions:} heart, kidney, and bone health are monitored
\end{itemize}

\section*{Complications}
\begin{itemize}
    \item Drug resistance from missed doses, limiting future options
    \item Side effects: nausea, insomnia, headache, and rarely serious toxicity
    \item Effects on kidneys, bones, and metabolism with long-term use
    \item Drug interactions with other medications and supplements
    \item Lipodystrophy and metabolic changes with older drugs, less common now
    \item Stigma and adherence difficulties
\end{itemize}

\section*{Prognosis and Outcomes}
Antiretroviral therapy is one of the great successes of modern medicine. People with HIV who start treatment promptly have a life expectancy close to that of the general population, a near-normal quality of life, and an undetectable viral load, meaning they cannot pass the virus to sexual partners or children. Treatment is lifelong, but regimens are simple and well tolerated, and monitoring prevents most long-term complications. With adherence, the prognosis for HIV is now excellent.

\section*{Prevention and Self-Care}
\begin{itemize}
    \item Start treatment as soon as possible after diagnosis
    \item Take medicines exactly as prescribed, at the same time daily
    \item Never stop treatment without medical advice
    \item Attend regular viral load, CD4, and side-effect monitoring
    \item Tell all health providers about HIV medicines to avoid interactions
    \item Check medicines before using over-the-counter or complementary products
    \item Maintain a healthy lifestyle: exercise, diet, and no smoking
    \item Use reminders or pill organisers to support adherence
\end{itemize}

\section*{Sources and Further Reading}
The following reputable sources provide further information for healthcare students and patients seeking reliable, current advice about HIV medicines.

\begin{itemize}
    \item ASHM. \textit{HIV antiretroviral therapy guidelines.} https://www.ashm.org.au/
    \item Healthdirect. \textit{HIV treatment and antiretroviral medicines.} https://www.healthdirect.gov.au/hiv-aids
    \item Positive Life NSW. \textit{Treatment information for people living with HIV.} https://www.positivelife.org.au/
    \item World Health Organization. \textit{Antiretroviral therapy.} https://www.who.int/
    \item NHS. \textit{HIV treatment: antiretroviral therapy.} https://www.nhs.uk/conditions/hiv-and-aids/
    \item HIV i-Base. \textit{Independent HIV treatment information.} https://i-base.info/
\end{itemize}
